\documentclass[11pt,a4paper]{article}
\usepackage[margin=22mm]{geometry}
\usepackage{amsmath,amssymb,amsthm,booktabs,longtable,array,fontspec}
\IfFontExistsTF{Noto Sans CJK SC}{\setmainfont{Noto Sans CJK SC}[AutoFakeBold=2,AutoFakeSlant=0.15]}{\setmainfont{Microsoft YaHei}[AutoFakeBold=2,AutoFakeSlant=0.15]}
\XeTeXlinebreaklocale "zh"
\XeTeXlinebreakskip=0pt plus 1pt
\setmonofont{lmmono10-regular.otf}[Scale=0.9]
\usepackage[colorlinks=true,linkcolor=blue,urlcolor=blue]{hyperref}
\usepackage{xurl,fancyvrb}
\DefineVerbatimEnvironment{verbatim}{Verbatim}{fontsize=\small}
\newcommand{\FS}{\operatorname{FS}}
\newcommand{\Z}{\mathbb Z}
\newcommand{\T}{\mathcal T}
\newcommand{\A}{A_*}
\newtheorem{theorem}{定理}[section]
\newtheorem{lemma}[theorem]{引理}
\renewcommand{\proofname}{证明}
\renewcommand{\abstractname}{摘要}
\renewcommand{\arraystretch}{1.14}
\setlength{\emergencystretch}{2em}
\linespread{1.07}
\raggedbottom
\title{独立闭合上界 84\\\large 尖锐源包络、三点截断与大小 85 的完整接入}
\author{Sum-of-Subset 研究项目}
\date{2026 年 10 月 19 日}
\begin{document}
\maketitle
\begin{abstract}
接续提交 \texttt{0a06472} 已独立核查的上界 85，本轮完成大小 85 的结构归约、
13 类实际情形与全部必要有限依赖的核查，得到 $71\le N\le84$。
扣除已计入中央区间的重新插入配对标签，将通用源参数界 $\Theta\le165$ 加强为
尖锐界 $\Theta\le162$，并给出实际 77 点源的等号例子。
精确空隙标签计数将全非空覆盖的空洞代价加强到 $w\ge2|D|+8$。
单空单元终端从 $107-w$ 加强到 $112-w$；最后一个零洞分支的 7315 状态枚举，
由两个一元不等式及每侧至多删除三点的实际子核心替代，统一矛盾余量至少 16。
新增验证器与完整归档套件均重新运行；程序只认证声明的有限命题，文字论证承担一般适用性。
\end{abstract}

\section{结论、先行上界与证据边界}
对有限实数集 $A$，令 $\FS(A)$ 包含空集和，$4A$ 允许重复元素。
记满足 $\FS(A)\subseteq4A$ 的最大可能基数为 $N$。
\begin{theorem}
接续已核查的上界 85，有
\[\boxed{71\le N\le84.}\]
\end{theorem}
\begin{proof}
只需排除大小恰为 85 的反例。若它不含零，加入零得到 86 点反例，违背前序上界。
为每个子集和选定一个四项表示，并指定零坐标，得到有限齐次整数线性方程组。
其有理核点在实核中稠密；选取足够接近原配置的有理核点保持全部严格次序，清分母后
得到 85 个互异整数。这是存在性归约，不假设原实数坐标本身共格，也不使用任意子集可行性。
第 \ref{sec:structure} 节的结构核查及第 \ref{sec:cases} 节的穷尽分类排除这个整数反例。
71 点构造保持为 $([-31,31]\cap\Z)\cup\{\pm93,\pm155,\pm217,\pm713\}$，
完整套件再次核对其 3349 个子集和值的四项表示。
\end{proof}

本报告与固定的归档原证明共同组成接入论证：
\begin{itemize}
\item 前序独立上界及一般基础：\path{reports/upper85_adoption_20261019.tex}。
\item 本次大小 85 的完整原证明：\path{verification/upper84/proof/report.tex}，
及 \path{verification/upper84/progress85/} 内的结构、逆定理、零洞、正洞和高洞证明。
\item 原四极值覆盖、固定标签和预算：
\path{verification/upper84/audit/verification_program/proof/report.tex}。
\end{itemize}
归档中“先行上界 86 的证明未随附件给出”的历史限制，现在由已发布的独立链补足。
归档字节保持不变，本报告明确替换其哪些步骤。此次独立范围与归档 84 一致；精确最大值仍未知。

\section{尖锐的源包络：从 165 到 162}\label{sec:source}
结构归约后 $|D|=77$，内核 $E=P_s\cup\{0\}\cup(-P_t)$ 有 69 点。
设 $r_s+r_t=68$，$R_v=\max P_v$，$R=R_s+R_t\ge68$。
每侧四个源尾点为
\[R_v+a_v, R_v+b_v, R_v+c_v, R_v+e_v,
\qquad 1\le a_v<b_v<c_v<e_v.\]
令 $K=|D+D|$、$U=a_s+b_s+a_t+b_t$、$C=c_s+c_t$、
$j_v=e_v-c_v$、$J=j_s+j_t$。与覆盖的空单元数 $j$ 区别，$j_v$ 只表示尾间隙。
定义整数源参数
\[\Theta=2K-4R-U-2C-J.\]
令 $\pi_v$ 为前三个偏移的不同二项和中大于 $c_v$ 的个数，
$B_v$ 为四个偏移的不同二项和中大于 $e_v$ 的个数，$u_v=a_v+b_v$。
已核对的四偏移引理及小三元组表给出
\begin{equation}\label{eq:fringe}
2\pi_v\le u_v+4-(u_v\bmod2),\quad
\pi_v+B_v-u_v\le5,\quad \pi_v\le u_v.
\end{equation}
最后一条在 $u_v\ge6$ 由 $\pi_v\le6$ 立即成立；$u_v=3,4,5$ 分别有
$\pi_v\le3,4,4$，且 $c_v\ge2b_v$ 后恒为 3。
第二条的无界证明保留：$u_v\ge9$ 用四元组界；$u_v\le8$ 分为
$c_v<2b_v$ 与恒定尾部，两个特殊三元组 $(1,3,4),(2,4,5)$ 在 $e_v\ge2c_v$ 后均有 $B_v=4$。

删去两侧最大值后，旧源有 75 点。先计中央整数区间
\[\mathcal I=[-2R_t-c_t,2R_s+c_s]\]
的全部 $2R+C+1$ 个位置，再计至多 $\pi_s+\pi_t$ 个例外。
重新插入两个最大值共增加至多 153 个形式配对标签。
若其中 $h$ 个已落在 $\mathcal I$ 内，它们不能再次增加这个上界。因此
\[
K\le2R+C+154+\pi_s+\pi_t-h,\qquad
K\le2R+C+J+1+B_s+B_t.
\]
中央区间无需已被旧二项和填满；这里从开始就计入了它的全部整数位置。
平均上述两界得
\begin{equation}\label{eq:average}
\Theta\le155-h+(\pi_s+B_s-u_s)+(\pi_t+B_t-u_t).
\end{equation}

\begin{lemma}\label{lem:162}
上述每个实际整数源都满足 $\Theta\le162$，且常数 162 可取到。
\end{lemma}
\begin{proof}
反设整数 $\Theta\ge163$。在删两最大值的界中先取 $h=0$，再用
\eqref{eq:fringe}，得到 $J\le316-\Theta\le153$。
考察七个不同的重新插入标签：正最大值与负侧三个旧尾点配对、反射的三个配对，
以及正负两最大值的配对。令
\[
n_s=\#\{x\in\{a_t,b_t,c_t\}:j_s\le R+x\},\qquad
n_t=\#\{x\in\{a_s,b_s,c_s\}:j_t\le R+x\}.
\]
前六标签在 $\mathcal I$ 内的个数恰为 $n_s+n_t$。
最后一个在区间内当且仅当
\[j_s-j_t\le R+c_t,\qquad j_t-j_s\le R+c_s.\]
若它不在区间内，不妨 $j_s\ge j_t+R+c_t+1$。
于是 $2j_t\le152-R-c_t\le81$，故 $j_t\le40<R+a_s$，反射的三个旧尾配对全部在内。
此时 $h\ge3$，\eqref{eq:average} 直接给 $\Theta\le162$。

因此可设最大值配对在区间内。若 $n_s+n_t\ge2$，仍有 $h\ge3$。
只剩下两种情形：
\begin{itemize}
\item $n_s=n_t=0$。这时 $j_s\ge R+c_t+1$、$j_t\ge R+c_s+1$，
故 $c_s+c_t\le15$。两侧均有 $j_v>c_v$，因而 $e_v>2c_v$、$B_v=4$。
由 $\pi_v\le u_v$，两项 fringe 超额之和至多 8；$h\ge1$ 给 $\Theta\le155-1+8=162$。
\item $n_s+n_t=1$。不妨 $n_s=1,n_t=0$。
于是 $j_s\ge R+b_t+1$、$j_t\ge R+c_s+1$，故 $c_s+b_t\le15$、$c_s\le13$。
又 $j_s\ge71>c_s$，所以 $B_s=4$。本侧超额至多 4，另一侧至多 5；
$h\ge2$ 给 $\Theta\le155-2+9=162$。
\end{itemize}
所有情形均矛盾。

等号取 $E=[-34,34]\cap\Z$，$D=([-37,37]\cap\Z)\cup\{-112,112\}$。
两侧偏移为 $(1,2,3,78)$，$R=68,U=6,C=6,J=150$，并且
\[D+D=([-149,149]\cap\Z)\cup\{-224,224\},\qquad K=301.\]
直接代入给 $\Theta=162$。这证明源包络的尖锐性，不声称这个源能扩充为可行的 85 点集合。
\end{proof}

记 $T_r=r(r+1)/2$，$\tau=\sum P_s+\sum P_t-T_{r_s}-T_{r_t}$。
实际覆盖中 $I=(r_s-34)^2$、$\A=I+\tau+\delta$，且
$\Theta=77j-280-w+\A$。因此统一加强为
\begin{equation}\label{eq:budget}
\boxed{I+\tau+\delta\le442-77j+w.}
\end{equation}
例如零洞 $Z_2,Z_3,Z_4,Z_5$ 的上预算分别为 $288,211,134,57$。
原各分类使用的较弱预算仍成立；原单最大值结论 $e_v\ge\max(4,\Theta-89)$ 保留，
在 $Z_5$ 中仍为 $e_v\ge\A+16$。

\section{全非空覆盖与单空单元的加强}\label{sec:gaps}
仍使用原 18 个普通整单元及两种配对，空单元指其全部实际目标贡献为空。
设 $M=\max D-\min D$，外侧大小 $v_1>v_2>v_3>v_4>M_v$，
$d_v=v_1-v_2,h_v=v_2-v_3,g_v=v_3-v_4,X_v=v_4-M_v$，$t_v=d_v-M$。
原中心记号 $O,H,\alpha,\beta,C4,C5,C7$ 不变。
两侧目标为 $F_v=\FS(D_v)\cap(0,\infty)$，其中 $D_s=D,D_t=-D$；
$S_v$ 为其最大值，$\eta_v=S_v-|F_v|$，$w=\eta_s+\eta_t$。
原松弛为 $c=2K+18|D|+85-|\FS(D)|$。

\subsection*{精确固定标签只计一次}
归档 \path{progress85/common/raw_gap_inventory.md} 的逐标签证明已核对。
第一空隙 $(H_v+M_v,O_v-M_o)$ 最多有 11 个固定标签：五个间距 $d_v$ 的链各取至多一项，
另有两个可同时出现的内部标签和四个跨侧标签。$C44,F5$ 不能同时出现的论证使用第一空隙的下端保护。
beta 活动的第二空隙最多有 8 个标签；两个空隙之并最多有 15 个。
后一计数按 $v_3<d_v$、$d_v\le v_3\le2d_v$、$v_3>2d_v$ 三段，内部至多 7 个，跨侧至多 8 个。
标签数不要求数值互异，重合只会降低覆盖能力。
在对应活动前提已保护空隙处，得到
\[
t_v\le12+\eta_v,\qquad
t_v\le9+\eta_v\ (\beta_v,H_v\text{ 活动}),\qquad
2t_v\le17+\eta_v\ (\beta_v,O_v\text{ 活动}).
\]
两个空隙共用同一批空洞与固定标签；不能把 15 改成两个独立预算。

\begin{lemma}\label{lem:noempty}
对整数源的原实际四极值覆盖，若 18 个普通单元全部贡献，则
\[w\ge2M+10\ge2|D|+8.\]
故 $|D|=77,w\le161$ 时至少一单元为空，原覆盖松弛 $c\ge77$。
\end{lemma}
\begin{proof}
四个同侧中心全活动，其整数线性组合恢复 $\lambda_v$ 和全部外侧点，
故 $X_v,g_v,h_v,d_v$ 为正整数。前序报告中由实际活动直接展开的三个不等式仍为
\[
\begin{aligned}
C5_v:&\quad t_v\ge2X_v+X_o+g_s+g_t+h_o+1,\\
C7_v:&\quad t_v\ge2X_v+X_o+g_v+1,\\
C4_v:&\quad t_v\ge2X_v+X_o+g_s+g_t+h_o+d_o+1.
\end{aligned}
\]
A 型中选实际 beta 贡献者 $s$；两侧 $C5$ 均活动。
恒等式 $C4_s+\beta_s+C5_t=-2s_3-2s_4-t_2-t_3-t_4$ 排除同侧 $C4_s$ 活动，
故第一对由 $t$ 贡献。于是 $t_s\ge7$，且 beta 与 $C4_t$ 给
$t_t\ge2M+2t_s+2X_s+2X_t+g_s+g_t+2\ge2M+22$。
第一空隙给 $\eta_t\ge2M+10$。
B 型两 beta 都活动；若第一对由 $s$ 贡献，$C7_t$ 给 $t_t\ge5$，
同样得到 $t_s\ge2M+18$，双空隙给 $\eta_s\ge4M+19\ge2M+10$。
最后 $M\ge|D|-1$，空整单元损失完整容量 $|D|$。证明无需核心分别本原或主导活动。
\end{proof}

\subsection*{同一个实际子核心的两种截断}
在原 $P_1/H_1$ 终端形状中，唯一空单元为一侧 $O$，A 型 beta 的唯一贡献者为另一侧 $u$；
外侧整性及共同条带已证明，$w$ 偶且 $\eta_u=w/2$。
侧大小为 $16$ 至 $52$，首缺失值 $m\le8$。保留低于 $m$ 的实际点，
其余从 $m+1$ 连续填充，记基线跨度 $B$、质量超额 $b$、实际额外质量 $\sigma$ 和径向超额 $\xi=R-B$。
保留的首洞证明给 $I+b-5B\ge-282$；$I+\tau\le395$ 排除 $m\ge9$，
低点每侧至多 3 个。高秩偏差递增，故 $\sigma\le52\xi$。

每侧删除最大的 $t=5$ 或 $6$ 点，得到同一个实际 $Q_t$ 服务于全部六条边，
\[
|Q_t|=69-2t,\qquad
f_t=137-2t-B+(\xi-\sigma)/t.
\]
原两目标的树 $L_v-V_v-U_v$ 及边 $L_v-C7_v$，仅以 $L_u$ 为根。
不重复计费的预算消元给
\[
2t_u\ge461-12t-5B+I+b+(1+6/t)\xi+(1-6/t)\sigma-w.
\]
$t=6$ 给 $107+2\xi-w$，$t=5$ 给 $119+(11\xi-\sigma)/5-w$。
当 $\xi\ge2$ 时前者至少 $111-w$；当 $\xi=0,1$ 时后者至少 $554/5-w>110-w$。
因为 $2t_u+w$ 为偶整数，所以
\begin{equation}\label{eq:112}
\boxed{112-w\le2t_u\le17+w/2,\qquad w\ge64.}
\end{equation}
原终端为 $107-w$ 及 $w\ge60$。这里先取得严格实数下界，再使用偶性，不能提前取整。

\section{最后一个零洞分支：以解析式替代 7315 状态}\label{sec:alpha}
考虑 $Z_5$ 的 A0 双空对、额外空 $\alpha_s$ 情形。
原实际活动、固定标签和正常端点论证继续沿用：先由核命中得到初始信用 14，
再通过 32 类上尾证明信用 27，随后才能证明端点正常。
进一步推出 $1\le C5_s<t_4$、$1\le\alpha_t,C5_t<s_4$，
LOW 标签增加后最终有 21 个 HIGH、47 个 LOW、触及 26 个原基线组，信用为 42。
这一步绝不提前假设端点正常；其完整标签与角色见
\path{progress85/foundations/z5_alpha_final85.md}。

写 $R_v=r_v+\xi_v$，$0\le\xi_v\le\tau_v$。
对每侧定义
\[
p_v=\min\left(\xi_v,\ \min_{1\le j\le3}\left\{j+\left\lfloor\frac{\tau_v-\xi_v}{j}\right\rfloor\right\}\right).
\]
将本侧点写成 $x_i=i+d_i$，其中 $0\le d_1\le\cdots\le d_{r_v}=\xi_v$。
删除最大的 $j\ge1$ 点时，$j d_{r_v-j}\le\tau_v-\xi_v$，所以保留侧的跨度至多
$r_v-j+\lfloor(\tau_v-\xi_v)/j\rfloor$。选择取到最小值的实际截断，
零截断表示保留全侧。合成一个含零的实际子核心 $Q$，对整数 $x\ge0$ 有
\[q=|Q|\ge63,\qquad |Q\cap(Q+x)|\ge F-x,\qquad F=69-p_s-p_t.\]
这是保持基数和跨度的同一次选择，不将不同截断的重叠相加。

令 $L(p)=\max(0,p-1,2p-4,3p-9)$。
由定义，$\xi_v\ge p_v$，$\tau_v-\xi_v\ge L(p_v)$。
因此有两个无界解析估计
\[
4\tau_s-2\xi_s-11p_s\ge4L(p_s)-9p_s\ge-21,\qquad
4\tau_t-11p_t\ge4L(p_t)-7p_t\ge-13.
\]
只需检查 $p=0,\ldots,5$ 的六行；当 $p\ge6$ 时 $L(p)=3p-9$，
两个表达式成为 $3p-36$、$5p-36$，均递增。
\begin{center}\small
\begin{tabular}{c|rrrrrr}
\toprule $p$&0&1&2&3&4&5\\\midrule
$L(p)$&0&0&1&2&4&6\\
$4L(p)-9p$&0&-9&-14&-19&-20&-21\\
$4L(p)-7p$&0&-7&-10&-13&-12&-11\\\bottomrule
\end{tabular}
\end{center}
所以不需要总质量上限或任何 7315 元状态清单，就有
\begin{equation}\label{eq:851}
4\tau-2\xi_s+2q+11F\ge2\cdot63+11\cdot69-21-13=851.
\end{equation}

原实际分组仍是：$\{\alpha_t,C5_t\}$、$\{H_s,U_s,V_s\}$、
两对 $\{U_t,V_t\},\{C6_s,C7_s\}$，以及未使用的 $C5_s,H_t$ 原生损失。
每组只占原保留整单元的容量；另用同一 $Q$ 的加权星形不等式作正线性组合，得到
\[
2M_t+2q+11F\le2R+827+2r_t-\delta+5(\delta-42).
\]
代入 $e_t\ge I+\tau+\delta+16$、$R=68+\xi_s+\xi_t$，再用新的
$I+\tau+\delta\le57$，将迫使
\[4I+4\tau-2\xi_s+2q+11F\le835,\]
与 \eqref{eq:851} 矛盾，余量至少 16。
若仅保留旧上预算 60，这个三点截断证明仍有余量 10。
原四点截断、$q\ge61$ 和 7315 状态的最小余量 6 不再是此分支的必要证明依赖。

\section{大小 85 的结构接入}\label{sec:structure}
完整推导见 \path{progress85/foundations/structure85_ap_bootstrap.md}、
\path{progress85/inverse/proof85_inverse.md} 及整合清单。
本轮核对的是实际集合的适用性与先后顺序，不是仅检查其标题或哈希。

\begin{enumerate}
\item 在固定 $|D|=77$ 前，缺失一侧的安全上界为 $24\cdot85+77=2117$，
小于单侧 70 点下界 2416。以后剥离中若一侧消失，下界 2347 大于上界 1925。
未配对数 $q$ 为偶且 $q^2\le2936$，故至少保留 7 对相反点。
先用 $\rho_s\ge-8,\rho_t\ge0$，得到小侧至少 10。
\item 标量反射只在此处取 $a=34-k,b=34+k$、$I=k^2$。
$\Lambda=1191+I,\Omega=600-I$。10、11、12 边界分别用 AP、严格逆定理和
有效洞见证 $(56,5,132)$；最大值删除所需的小侧 $+16$ 修正保持。
12 边界最后是 $\rho_Q\ge125>124$。
\item 大侧 AP 由剩余类与无裁剪空隙给共同格。仅小侧 AP 的两剩余类先证明其阶为 2，
再用实际奇偶生成元质量推出最小奇生成元 $\le15$，从而 $z\le14$；
没有在质量论证之前使用 $h_Q\le\kappa$。
整性和两层相交后，非本原与半步长分别给 $c\le59,61<77$，$w\le56$，
由引理 \ref{lem:noempty} 统一排除。AP 的 13、14 边界同样闭合。
\item 双非 AP 的充分严格逆检验只有 $k=0,1,2$ 自动成立。
实际失败的 $k=3,\ldots,20$ 用 18 个新端点修复；其中 $k=10,13,14$ 必须使用
非 AP 前缀的加强起点，不能沿用较弱的 14 阶起点。
严格排离格和端点桥后，通过短直径定理得到精确 $h_P,h_Q\le\kappa$，
不提前假设另一侧本原。
\item $a=13$ 的失败分支用 $(55,2,52)$ 得 $h_Q\le4$，最终 $185>159$。
成功分支先以 $\kappa\le110$ 和带修正的 $-811<0$ 完成格化，
然后由 $z\ge17$ 用 $(55,4,119)$ 得 $\rho_Q\ge113$，最终 $160>159$。
两分支必须分别处理。
\item 未修复严格逆分支只有 $z\ge83+I$ 需进一步证明末步相交。
$k\ge16$ 先排除；余下 32 个定向分支有 18 个直接矛盾，14 个用加强前缀端点修复。
得到实际末步相交后，最大值删除留下至少一半的洞，内层比较变为 $-165-2I<0$。
外层比较用 $1191>4\cdot77-10+600$。
\item 此后才使用整数源的 $G=144$、$w\le4z$ 与
\[c+\rho+\T\le600-I,\qquad \T\ge\max(48,z,144-2z,\kappa+144-3z).\]
修复分支若小侧非本原，$c\le28,w\le64$，由新全非空引理排除。
14 边界的严格分支用 $(54,6,168)$，或 $c\le31,w\le48$，完整排除。
\end{enumerate}
最终确立 $15\le r_s,r_t\le53$、两侧各自本原、共同格和两层实际相交。
正核心洞时有 $\mu\ge62,\tau\ge62+2z$。
活动分类中物理符号不再因侧大小重新命名；所有者必须是实际贡献者。

\section{全部情形及进一步简化}\label{sec:cases}
令 $j$ 为实际空普通整单元数，$\delta=c-77j\ge0$。
零核心洞 $z=0$ 分 $Z_0,\ldots,Z_5$；正低洞 $1\le z\le14$ 分 $P_0,\ldots,P_5$；
$z\ge15$ 为 $H$。更大的低洞空单元数由预算排除。
完整原论证按上文列出的 13 类逐一核对；以下索引记录具体接入与替换。

\begin{longtable}{p{.08\linewidth}p{.85\linewidth}}
\toprule 情形&本轮核对与收尾\\\midrule\endfirsthead
\toprule 情形&本轮核对与收尾\\\midrule\endhead
$Z_0,P_0$&共同整数源成立后，$w\le56<162$，引理 \ref{lem:noempty} 排除 $j=0$。\\
$Z_1$&保留原一空单元整性、唯一空 O、对侧 beta 形状及同一子核心六边森林；
下界 $2t_u\ge73+I$ 与新双空隙 $2t_u\le17$ 矛盾。\\
$Z_2$&核对双上帽、HIGH $C7$ 原生损失和四张 beta/O 卡，保留三种真实缺失形式。
加强后的最终六边矛盾见下文，余量至少 56。\\
$Z_3$&A0 核对五个加权端点的全部轻重位置；唯一重 alpha 的 11 个 LOW 标签触及 9 组，
与 8 个 HIGH 合为信用 10，紧邻终端 $792>791$。
A1 用物理侧 7 单元与主导的实际并集，四个端点位置上界最大 338，而下界 394。\\
$Z_4$&核对 A-E/A-H/B 的完整保留表、F/O/S/E 各卡的端点归属，
以及两个 beta 都缺失时的额外排除 $1192>1063$。
S 卡保留修正后的 $-3\delta/4$。
双空对不指定所有者，原信用 18、删除两最大值的 67 点源排除四种 H 核命中状态。\\
$Z_5$&先核命中，再外侧整性和完整 A0/A1/B 分类。
A0 十三张保留第一对卡、双空对 H/C5 卡、alpha 卡均接入；alpha 卡改用第 \ref{sec:alpha} 节解析证法。
分裂 beta 始终只计整个配对容量。
A1 最后 $C5_t$ 先证明 beta 截断，再由 32 类尾型删至 26、22 类，得到
$\tau+\delta\ge62$，与新上界 $57-I$ 矛盾。\\
$P_1$&核对四个内层实际发生项的全部子集位移及共同条带、双上帽排除和唯一 O 形状。
$w\le14$ 与新终端 $w\ge64$ 矛盾。\\
$P_2$&先完成共同条带、整性和首洞分类；利用尖锐空隙简化 $C7$ 及 beta/O 两处收尾，
最后得到 $100+I+\xi\le3w/2\le21$，见下文。\\
$P_3$&完整重算 $\binom{17}{12}=6188$ 个种子并核对无限传播；从 $z\le10$ 降到 8，
首洞 $m\le3$ 时再降到 6。整性、A0/A1 穷尽及 $m=1,2,3,4$ 全部端点收尾接入。
特别核对 clip 72、最多一个重端点、物理唯一 1 的两个所有者方向，保留修正后的全部常数。\\
$P_4$&七个种子先给六种洞型；信用 8 排除四种高质量洞型，仅留首洞 1、2。
整性、CORE8、A0/A1/B 所有者规则均在正常端点之前确立。
接入对侧／同侧所有者、T5 和 A0 三卡；首洞 1 的交叉卡现由 $\Theta\le162$ 直接给余量 6，
不再讨论端点 164、165 的 fringe 修补。\\
$P_5$&共同条带给 $w\le z$。源预算变为 $I+\tau+\delta\le57+w$，
与 $\tau\ge62+2z$ 直接矛盾，无需外侧整性。\\
$H$&配对的 $(\mu,\rho)$ 两个下界都保留，先完成整个高洞域的共同条带。
$j\ge2$ 源预算排除；$j=0$ 用新全非空引理；$j=1$ 只余 $z=16,18,20$，由新终端排除。\\
\bottomrule
\end{longtable}

\subsection*{零洞两空单元的强终端}
保留原分类后，选取实现 $f=\max(69-\xi,57-(\tau-\xi)/6)$ 的同一实际子核心，服务于两侧六边。
至少一侧有上帽，精确空隙给 $t_u\le12$；实际 $\alpha_o,C5_u$ 给
$t_u\ge2X+G+2$，故 $2X+G\le10$。于是
\[
\A\ge I+\tau+6f-10\ge I+\max(414-5\xi,342+\xi)-10\ge344+I.
\]
两个仿射式在 $\xi=12$ 相交，最小值 354。新源预算 $\A\le288$ 给至少 56 的矛盾余量；
旧收尾的余量仅为 2。

\subsection*{正洞两空单元的短路径}
保留原条带与整性证明，并暂用较弱但有效的 $\A\le291+w\le305$。
若 $C7_v$ 为空，原 HIGH 损失给 $\delta\ge4r_v+96$。
首洞至少 2 时 $\tau\ge85$，故 $\A\ge313$ 已矛盾。
首洞 1 时写 $\tau=68+\sigma$，得到 $(r_v-32)^2+\sigma\le w-5\le9$。
两侧各含连续前缀 $2,\ldots,21$，其余核心点至多 49；种子质量 230 的双洞支持持续。
于是 $z=w=2$，但右侧 $w-5<0$，仍矛盾。原四个额外固定信用及最后 fringe 等号分析均可省去。

之后沿原根森林把首洞降到 1，cap 侧保留 V。
若 beta 和 O 同侧活动，根森林与新双空隙给
$82+4I+\sigma\le5w\le70$，直接矛盾，不再单独处理 $w=14$。
同样只剩三种实际缺失形式，$2X+G\le10+w/2$。
六点截断给 $2X+G\ge330-\sigma+\xi-\delta$，而新源预算为
$I+\sigma+\delta\le220+w$，合并得到 $100+I+\xi\le3w/2\le21$。

\subsection*{高洞先局部化，再用新引理收尾}
在四段 $z=15{:}16,17{:}20,21{:}48,\ge49$，原配对下界依次为
\[(\mu,\rho)\ge(426,294),(470,338),(514,382),(554,422).\]
必须保留 $\mu$ 信息，不能仅由 $\rho\le\mu$ 倒推出它。
核心中央宽度减最大源系数的余量分别至少 $195,315,363$，尾段至少 $93+7I+6c$；
两侧质量超过边界空洞数的余量为 $171,199,131,120$。
故同一条带传播过全部八个源尾点，得 $w\le z$、偶性与 $\eta_v=w/2$。
这些传播条件仅使用局部化之前的预算，没有循环。

$j\ge2$ 时 $I+\tau+\delta\le288+w$ 与 $\tau\ge\rho+2z$ 矛盾。
$j=0$ 时新源预算给 $I+c+\rho+z\le442$：尾段至少 $422+49>442$；
其余 $w\le48<162$，由引理 \ref{lem:noempty} 排除。
$j=1$ 时原实际双上帽论证留下唯一 O、对侧 beta 的终端形状，
源预算仅允许偶数 $z=16,18,20$，由 \eqref{eq:112} 的 $w\ge64$ 排除。

\section{有限验证、复现与交付}\label{sec:verification}
新的入口固定 106 个源文件，复制到全新临时目录，执行归档完整套件与新增检查，
结束时再次核对源字节、11 个顶层套件、17 个新模块（含 7 个聚合子模块），
及继承的 140 个唯一检查编号。运行记录作为输出，不作为通过依据。
归档中的一般定理、源到模型的实际蕴涵和无界归纳均由文字证明承担。

两次完整回放均为 11/11 顶层套件通过，继承检查均为 140 PASS、0 FAIL、0 SKIP。
去除环境与计时字段后，逐项数学记录一致；新增几何结果逐字节一致。
新增程序检查三类中央标签扣除的算术证书、七标签的真实区间归属、
28,900 个实际源的回归族及等号集合；该有限族不充当无界包络的证明。
首洞枚举有 35 个低前缀模式、56,662 个必要标量状态，取偶数前最小值为 111。
三点截断的不等式以六行及正斜率尾部给出无界证明，另核对 80,601 个单侧参数作为回归。
7 项故意抬高的证书界或删改的源清单全部被拒绝。旧 7315 状态检查仍完整重放作为额外交叉检查。

在仓库根目录运行：
\begin{verbatim}
python -I -S -B verification/upper84_adoption.py --output .work/upper84.json
python -I -S -B -O verification/upper84_adoption.py --output .work/upper84_O.json
python -I -S -B experiments/check_upper84_controls.py --output .work/controls84.json
\end{verbatim}
仅需要 Python 3.10+ 标准库。\texttt{-O} 启动时，旧模块仍强制保留断言；新增检查使用显式判定。
Windows 启动器为 \path{verify_upper84_windows.bat}。
本轮实际运行环境为 Linux，另检查带空格和中文的隔离路径；不声称实际运行过 Windows。
本报告使用 XeLaTeX，优先 Noto Sans CJK SC，缺失时回退 Microsoft YaHei。

实际运行回执、逐项覆盖、哈希、完整日志与 PDF 审阅记录位于
\path{results/2026-10-19/upper84/}。
有限 PASS 不等于证明助手的形式化认证。
下一步是研究大小 84 的排除；仅调整这里的常数不能自动得到上界 83。
\end{document}
